\begin{proof}
We use the extended-real convention stated before
\Cref{thm:ch16-collapse-finite-tropical-curves}.  First note that
\(\mathcal J_X\) is a sheaf.  Only the junction condition needs checking.
Suppose it holds on an open cover and consider one displayed junction tuple.
An open member containing \((0,y_0)\) contains, for all sufficiently small
\(s>0\), the segments from \((0,y_0)\) to
\[
 \bigl(st_av_{i(a)},\,y_0+s(y_a-y_0)\bigr).
\]
Applying the local junction inequality to these contracted points and then
using branchwise convexity gives the required inequality after cancelling
\((1-s)u(0,y_0)\) and dividing by \(s\).  If the central value is
\(-\infty\), there is nothing to prove.  Joint extended-real convexity and
the junction inequalities are both preserved by pointwise decreasing
limits.  Finally, a generator belongs to \(\mathcal J_X\): its ambient
tangent affine function at \((0,y_0)\) is a branchwise minorant, and its
weighted linear terms cancel at the displayed barycentre.  Hence
\[
                         \mathcal P_X\subseteq\mathcal J_X.
 \tag{\(\dagger\)}
\]

It remains to prove the reverse inclusion at the first stage.  Away from
\(\{0\}\times\Omega\), this is the one-point case of
\Cref{thm:p071-stone-euclidean-collapse}, after choosing a ray coordinate
and shrinking to an ordinary convex ball.  We therefore work at
\((0,y_0)\), with \(u\in\mathcal J_X\) defined on a product neighbourhood.
Shrink once so that this neighbourhood contains
\[
\bigl(\{0\}\cup\coprod_i(0,2T_i]v_i\bigr)\times B_0,
\qquad 0<2T_i<r_i,
\]
where \(B_0\Subset\Omega\) is an ordinary open convex ball about \(y_0\).
Put \(c(y):=u(0,y)\).
For every branch write \(u_i(0,y):=c(y)\) and
\(u_i(t,y):=u(tv_i,y)\) for \(t>0\); these are jointly convex on their
closed branch half-chambers by the defining comparison condition.

If \(c(y_0)=-\infty\), then after a further product shrinking \(u\) is
identically \(-\infty\).  Indeed, every point \((tv_i,y)\) in a sufficiently
small product is the midpoint of \((0,y_0)\) and
\((2tv_i,2y-y_0)\), both lying in the original product.  Branchwise
convexity gives value \(-\infty\) at that midpoint.  The constants \(-n\)
then give the required approximation.  We may consequently suppose, after
shrinking the ball, that \(c\) is finite on \(B_0\): if \(c\) were
\(-\infty\) at a nearby point, expressing \(y_0\) as a strict convex
combination of that point and a point slightly beyond \(y_0\) would
contradict \(c(y_0)\in\mathbf R\).  A branch is called
\emph{polar} if its restriction to
\((0,T_i)\times B_0\) is identically \(-\infty\).  Every other branch is
finite there: if a finite interior point and a \(-\infty\)-valued interior
point existed,
extend the line through them a little beyond the finite point inside the
open convex chamber and use convexity to contradict finiteness.  Moreover,
no nonnegative relation
\[
                         \sum_i\lambda_iv_i=0
 \tag{\(\ddagger\)}
\]
uses a polar index with positive coefficient.  Otherwise take all Euclidean
coordinates equal and all normal parameters equal and small in the junction
inequality; its right-hand side is \(-\infty\), contradicting the finite
central value.

Choose numbers \(0<\rho_i<R_i<T_i\), a compact convex set
\(Q\Subset B_0\), and an ordinary open convex ball \(B\) with
\(\overline B\subset\operatorname{int}(Q)\) and \(y_0\in B\).  On every
finite branch, \(u_i\) is bounded above and below on
\([0,T_i]\times Q\).  For the upper bound interpolate along the normal
coordinate between \(t=0\) and one fixed positive level in
\((T_i,2T_i)\).  For the lower bound choose two positive levels in
\((T_i,2T_i)\), express the
first as a convex combination of \(t\in[0,T_i]\) and the second, and solve
the resulting convexity inequality for \(u_i(t,y)\).  The positive-level
slices and \(c\) are finite convex functions on a neighbourhood of \(Q\),
hence are bounded there.

For a finite branch define, for \(p\in(\mathbf R^d)^*\),
\[
 f_i(t,p):=\min_{y\in Q}\bigl(u_i(t,y)-p\mathbin\cdot y\bigr),
 \qquad
 f_0(p):=\min_{y\in Q}\bigl(c(y)-p\mathbin\cdot y\bigr).
 \tag{\(\star\)}
\]
The minima are finite, and \(t\mapsto f_i(t,p)\) is convex.  For the
latter assertion, choose minimizing Euclidean variables at two normal
parameters, take their convex combination in \(Q\), and apply joint
convexity.  On a polar branch put \(f_i(t,p)=-\infty\) for \(t>0\), with
the common central value \(f_0(p)\).

For each fixed \(p\), the resulting star profile satisfies the finite-star
junction inequalities of \Cref{thm:ch16-collapse-finite-tropical-curves}.
Indeed, for a normal junction tuple choose arbitrary \(y_a\in Q\), put
\(\bar y=\sum_a\alpha_ay_a\), and use the product junction condition to get
\[
 \begin{split}
 f_0(p)
 &\leq c(\bar y)-p\mathbin\cdot\bar y\\
 &\leq\sum_a\alpha_a
       \bigl(u_{i(a)}(t_a,y_a)-p\mathbin\cdot y_a\bigr).
 \end{split}
 \tag{\(\star\star\)}
\]
Taking the independent infima proves the claim.  Conversely, if profiles
\(h_i(t,p)\) satisfy those finite-star inequalities, then
\[
 H_i(t,y):=\sup_p\bigl(p\mathbin\cdot y+h_i(t,p)\bigr)
 \tag{\(\clubsuit\)}
\]
satisfies the product junction condition: apply the finite-star inequality
at each \(p\), use the Euclidean barycentre identity, and then take the
supremum.  Each expression in \((\clubsuit)\) is jointly convex in
\((t,y)\), so the same is true of \(H_i\).

For each finite branch and \(y\in\operatorname{int}(Q)\), the elementary
local biconjugacy identity is
\[
 u_i(t,y)=\sup_p\bigl(p\mathbin\cdot y+f_i(t,p)\bigr).
 \tag{\(\heartsuit\)}
\]
The inequality from right to left follows from \((\star)\).  For the other
one, take a supporting covector of the finite convex function
\(u_i(t,\cdot)\) at \(y\); it realizes equality in \((\star)\).  The same
argument applies to \(c\) at \(t=0\); on polar branches both sides are
\(-\infty\).

Choose nested finite partitions of every \([0,T_i]\), with mesh tending to
zero and with both endpoints included.  They are chosen independently of
\(p\).  On a finite branch let \(h_{n,i}(t,p)\) be the chordal interpolant
of \(f_i(\cdot,p)\) on the \(n\)-th partition, and on a polar branch set
\[
                         h_{n,i}(t,p)=f_0(p)-nt.
\]
By convexity, the finite-branch chordal interpolants decrease to \(f_i\),
and they majorize it.  Relation \((\ddagger)\) and
\((\star\star)\) therefore show that every profile \(h_n(\cdot,p)\)
satisfies the finite-star junction inequalities.  Define \(H_{n,i}\) from
\((\clubsuit)\) with these profiles.  Then \(H_{n+1}\leq H_n\), and on a
finite branch, if \(t\) lies between two adjacent partition points \(a,b\)
with affine weights \(1-\theta,\theta\), \((\heartsuit)\) gives
\[
 u_i(t,y)\leq H_{n,i}(t,y)
 \leq(1-\theta)u_i(a,y)+\theta u_i(b,y).
 \tag{\(\spadesuit\)}
\]
Consequently \(H_n\downarrow u\) pointwise on the smaller product.  On a
polar branch the formula is simply \(H_{n,i}(t,y)=c(y)-nt\) for
\(y\in\operatorname{int}(Q)\).

For fixed \(n\), \(H_n\) is continuous on the compact core
\[
 K:=\Bigl(\{0\}\cup\coprod_i(0,R_i]v_i\Bigr)\times\overline B.
\]
It is finite convex, and hence continuous, in every open chamber.  At the
common face, the first chord interval in \((\spadesuit)\) gives the required
upper limit, while the supremum in \((\clubsuit)\) gives lower
semicontinuity and \(H_n(0,y)=c(y)\).  The same chord estimate handles the
partition knots and the outer compact boundary.  Finiteness of the star then
gives continuity on all of \(K\).

Fix a summable positive sequence \((\varepsilon_n)\).  Compactness of \(K\)
and continuity just proved give, for every \(n\), finitely many covectors
\(p_{n,1},\ldots,p_{n,q_n}\) such that
\[
 H_n-\varepsilon_n
 \leq A_n:=\max_j\bigl(p_{n,j}\mathbin\cdot y+
                    h_{n,i}(t,p_{n,j})\bigr)
 \leq H_n
 \quad\text{on }K.
 \tag{\(\diamondsuit\)}
\]
Here the middle expression is read on branch \(i\).  To obtain this finite
set, choose one near-maximizing covector at each point of \(K\) and pass to
a finite subcover; no compactness of the full dual space is asserted.

For each of these finitely many scalar star profiles, repeat the quantitative
common-affine-minorant and regularized-maximum construction in the proof of
\Cref{thm:ch16-collapse-finite-tropical-curves}, choosing its error at most
\(\varepsilon_n\).  It yields smooth branchwise-convex functions
\(g_{n,j}\) on the fixed smaller substar such that
\[
 h_n(\cdot,p_{n,j})\leq g_{n,j}
 \leq h_n(\cdot,p_{n,j})+\varepsilon_n.
\]
The functions
\[
               (tv_i,y)\longmapsto g_{n,j}(tv_i)+p_{n,j}\mathbin\cdot y
\]
are ambient-smooth and jointly convex on every branch.  A convex,
coordinatewise nondecreasing regularized maximum of this finite family is
again of that form.  Using regularization error \(\varepsilon_n\) and then
adding one further \(\varepsilon_n\), we obtain
\(w_n\in\mathcal G_{X,0}\) on the fixed product
\[
 V:=\Bigl(\{0\}\cup\coprod_i(0,\rho_i)v_i\Bigr)\times B
\]
with
\[
                         H_n\leq w_n\leq H_n+3\varepsilon_n
 \quad\text{on }V.
 \tag{\(\triangle\)}
\]
The ambient smoothness at the common face is legitimate because only
finitely many scalar profiles occur for a fixed \(n\); their affine collars
therefore have one common positive radius.  That radius may depend on \(n\)
without changing \(V\).

Finally set
\[
 d_n:=3\sum_{k>n}\varepsilon_k,
 \qquad
 G_n:=w_n+d_n.
\]
The inequalities \((\triangle)\) and \(H_{n+1}\leq H_n\) give
\[
 G_{n+1}\leq H_{n+1}+3\varepsilon_{n+1}+d_{n+1}
 \leq H_n+d_n\leq G_n.
\]
Since \(H_n\downarrow u\) and \(3\varepsilon_n+d_n\to0\), the sequence
\(G_n\) decreases pointwise to \(u\) on the one fixed neighbourhood \(V\).
Thus \(\mathcal J_X\subseteq\mathcal P_{X,1}\).  Combined with
\((\dagger)\) and the tautological inclusion
\(\mathcal P_{X,1}\subseteq\mathcal P_X\), this proves the theorem.
\end{proof}
